\lecture{3}
\title{Frames, Monogram Notation, and Rotations}

\begin{document}

\subsection{Frames}

Note that $\mathbb{E}^3$ denotes an \emph{affine space}.
Informally, it's $\mathbb{R}^3$ without an origin---so
it has subtraction and distances, but no addition.
In robotics, we work in $\mathbb{E}^3$ instead of $\mathbb{R}^3$
because we oftentimes can't pick a single universally-meaningful origin.

\textbf{Definition.}
In robotics, a \textbf{frame} (similar to a \emph{basis} in linear algebra)
describes a frame of reference using the following data.
\[ \text{frame} ~ = ~
(\text{origin } O \in \mathbb{E}^3, ~
\text{basis } (\hat{x}, \hat{y}, \hat{z})),
~~~ \text{ where } 
(\hat{x}, \hat{y}, \hat{z})
\text{ must be orthonormal and right-handed.} \]


Note that every body $i$ has a frame $B_i$ associated with it.
(Imagine an airplane, for example.)

\textbf{Definition.} Drake has a built-in \textbf{world frame} denoted $W$.

\textbf{Definition.}
Any frame is defined by its \textbf{position} $O$
and its \textbf{orientation} $(\hat{x}, \hat{y}, \hat{z})$.
This data is collectively called \textbf{spatial pose}.

When we refer to ``the pose of an object $O$'',
we really mean ``the pose of the body frame of $O$ in the world frame $W$''.

\subsection{Monogram Notation}

\textbf{Definition.}
We write $^B p^A_C$ to denote the vector $\vec{BA}$ expressed in frame $C$.
(So $^B p^A_C$ is a tuple of $3$ real numbers.)  Furthermore,
\begin{itemize}
    \item Since all frames $F$ have a position,
    it is valid to write something like $^{F_1}p^{F_2}_{F_3}$.
    \item For any frame $F$, the notation $^Fp^A$
    is shorthand for $^Fp^A_F$.
    \item For a point $B$ that isn't a frame, $^Bp^A$ is shorthand for $^Bp^A_W$.
\end{itemize}

\textbf{Definition.}
For any two frames $F$ and $G$, we say:
\begin{itemize}
    \item The notation $^FR^G$ denotes
    the \textbf{rotation} that sends the orientation of $F$
    to the orientation of $G$.
    \item The notation $^FX^G$ denotes
    the \textbf{spatial transform} that sends the spatial pose of $F$
    to the spatial pose of $G$.
\end{itemize}

\textbf{Definition.}
The Drake documentation suggests the following variable names:
\[ ^Bp^A_C = \verb|p_BA_C|
~~~~~~~~~
^FR^G = \verb|R_FG|
~~~~~~~~~
^FX^G = \verb|X_FG| \]

\begin{theorem}{Monogram Algebra}
    The following identities are all true.
    \[
    {}^Ap^B_F + {}^Bp^C_F = {}^Ap^C_F
    ~~~ \parallel ~~~
    {}^Ap^B_F = - {}^Bp^A_F
    ~~~ \parallel ~~~
    {}^Ap^B_G = {}^GR^F {}^Ap^B_F
    ~~~ \parallel ~~~
    {}^AR^B {}^BR^C = {}^AR^C
    \]
\end{theorem}

\begin{proof}
    The first, second, and fourth of these are obvious.
    The third is surprisingly unintuitive---think through it carefully.
    ~ $\blacksquare$
\end{proof}

\begin{remark}
    There are two ways to think about the rotation $^AR^B$.
    \begin{itemize}
        \item It is a map from the orientation of $A$
        to the orientation of $B$.
        \item It is a map that, for any point $Q$,
        sends the coordinates of $Q$ in the frame of $B$
        to the coordinates of $Q$ in the frame of $A$.
    \end{itemize}
\end{remark}

We now dedicate the rest of this lecture
to giving examples of how we actually describe rotations numerically.

\subsection{Rotation Matrices}

Just a matrix $M \in SO_3(\mathbb{R})$.
The main nontrivial point to be aware of
is that $M^{-1}$ is very easy to compute---it's just $M^{-1} = M^{\top}$.

Note, also, that using $9$ numbers for a data structure
that describes rotation in $\mathbb{R}^3$ is, in some sense,
wasteful.  There are ways to describe rotations
using fewer than $9$ numbers, such as\dots

\subsection{Roll, Pitch, Yaw (Euler Angles)}

\textbf{Definition.}
Let $R_x(\theta)$ denote the rotation about the world's $x$-axis by $\theta$.
Define $R_y(\theta)$ and $R_z(\theta)$ similarly.

\begin{remark}
    Use the right-hand rule.
    If you point your thumb in the direction of the positive $x$-axis,
    then your fingers curl in the direction of positive rotations $\theta$.
\end{remark}

\textbf{Definition.}
Let $(r, p, y)$ in \textbf{roll-pitch-yaw} denote the rotation
$R_z(y)R_y(p)R_x(r)$ (read from right-to-left, as with matrix multiplication).
These are also called \textbf{Euler angles}.

\begin{center}
    \frame{\includegraphics[width=0.5\textwidth]{lecture-03/roll-pitch-yaw.png}}
\end{center}

\textbf{Example.}
The rotation with roll $90^{\circ}$, pitch $0^{\circ}$, and yaw $90^{\circ}$
has the effect of permuting the axes $(x, y, z) \mapsto (y, z, x)$.

\begin{theorem}{Extrinsic vs. Intrinsic}
    Given rotation angles $(r, p, y)$
    and an object $O$, consider the following rotations on $O$, where
    the \textbf{extrinsic} axes are the world's axes,
    whereas the \textbf{intrinsic} axes are the basis of the frame of $O$.
    \begin{itemize}
        \item Rotate by $r$ about the \textbf{extrinsic} $x$-axis,
        then rotate by $p$ about the \textbf{extrinsic} $y$-axis,
        then rotate by $y$ about the \textbf{extrinsic} $z$-axis.
        \item Rotate by $y$ about the \textbf{intrinsic} $z$-axis,
        then rotate by $p$ about the \textbf{intrinsic} $y$-axis,
        then rotate by $r$ about the \textbf{intrinsic} $x$-axis.
    \end{itemize}
    Then these two rotations are the same.
\end{theorem}

\begin{proof}
    Say the rotations about the \textbf{extrinsic} axes
    have rotation matrices $R_X$, $R_Y$, and $R_Z$.
    Then the intrinsic rotations are\dots
    \begin{itemize}
        \item ``rotate by $y$ about the intrinsic $z$-axis''
        $\implies$ this is equivalent to $R_Z$.
        \item ``then rotate by $p$ about the intrinsic $y$-axis''
        $\implies$ this is equivalent to $R_ZR_YR_Z^{-1}$ \dots
        \begin{itemize}
            \item \dots because it's equivalent to
            ``undo the previous rotation,
            then apply $R_Y$, then redo the previous rotation''.
        \end{itemize}
        \item ``then rotate by $r$ about the intrinsic $x$-axis''
        $\implies$ this is equivalent to
        $(R_ZR_YR_Z^{-1}R_Z)R_X(R_ZR_YR_Z^{-1}R_Z)^{-1}$ \dots
        \begin{itemize}
            \item \dots for the same reason as the previous bullet.
        \end{itemize}
    \end{itemize}
    And composing all three gives
    $(R_ZR_Y)R_X(R_ZR_Y)^{-1} \cdot R_ZR_YR_Z^{-1} \cdot R_Z = R_ZR_YR_X$,
    as desired. ~ $\blacksquare$
\end{proof}

The problem with using Euler angles
to parametrize the space of rotations
is that this parametrization produces \textbf{gimbal lock}.

\textbf{Example.}
Consider the subspace of roll-pitch-yaw rotations
$\{(r, 90^{\circ}, y) \mid r, y \in \mathbb{R}\}$.
This happens to be a \emph{one-dimensional} subspace,
despite there being $2$ free parameters $(r, y)$.

The following (warning: AI-generated) widget
gives a visual demonstration of gimbal lock when $p = 90^{\circ}$.

\includewidget{lecture-03/gimbal-lock}

\begin{remark}
    The ``only'' times when gimbal lock occurs
    are when $p = \pm 90^{\circ}$---you should
    intuitively be able to see why using the widget above.
    (What is the \emph{only} way that
    two axes can overlap with each other?)
\end{remark}

In fact, it has been proven (Stuelpnagel, 1964)
that no parametrization of $SO_3(\mathbb{R})$ with three variables
can avoid gimbal lock.

\subsection{Axis-Angle}

Recall from \href{/18.701/lectures/9/#orthogonal-groups}{18.701}
that every rotation in $SO_3(\mathbb{R})$ looks like
a rotation about some axis $\hat{v}$ through the origin
by some angle $\theta$.
So one natural way to encode a rotation
is through the three-variable vector
$\theta \cdot \frac{\hat{v}}{\|\hat{v}\|}$.

Of course, this parametrization also produces gimbal lock,
because it has exactly three variables.

\subsection{Quaternions}

See \href{/18.701/lectures/20}{18.701} for
a more rigorous discussion of the quaternions $\mathbb{H}$
(along with the group of unit quaternions $SU_2$).
Here's the main point we need to know.

\begin{theorem}{Unit Quaternions as Rotations}
    The \emph{unit} quaternion $(w, x, y, z) \in SU_2$
    represents a rotation about axis $\hat{v}$ by angle $\theta$
    (with $\|\hat{v}\| = 1$)
    if we have $w = \cos\left(\frac{\theta}{2}\right)$
    and $\begin{bsmallmatrix} x \\ y \\ z \end{bsmallmatrix}
    = \sin\left(\frac{\theta}{2}\right) \hat{v}$.
\end{theorem}

\textbf{Example.}
A rotation by $\theta = 60^{\circ}$ about the axis
$\hat{v} = \begin{bsmallmatrix} 2/3 \\ 2/3 \\ 1/3 \end{bsmallmatrix}$
can be expressed by the quaternion
$\left(\frac{\sqrt{3}}{2}, \frac{1}{3}, \frac{1}{3}, \frac{1}{6}\right)$.

\end{document}